\section*{Phase 1: Business Understanding}

%\textbf{Project Scope Document}

\subsection*{Task 1.1 Determine Business Objectives}

\subsubsection*{Background / Problem Management Wants to Address}

Contagious diseases, like influenza, result in a substantial public health challenge due to elevated hospitalization rates and increased spending along with the risk of death. The erratic nature of these outbreaks makes properly swelling them a challenge. This project seeks to accomplish the creation of a predictive model for outbreaks of influenza using previously gathered (\textcite{who_gisrs}) data.

\subsubsection*{Business Goals}

The primary goal of this project is to predict early detection and forecasting of influenza outbreaks to help, healthcare providers, and public health agencies to be prepared. Specific goals include:
\begin{itemize}
  \item Discovering and documenting various trends during influenzas breakout seasons.
  \item Formulating predictive algorithms for expected cases in the weeks to come.
  \item Improving public awareness and mitigation efforts through early warnings.
\end{itemize}

\subsubsection*{Business Success Criteria}

The success of this project will be measured using:
\begin{itemize}
  \item Accuracy of Predictions: The predictive model should achieve at least 80\% accuracy in forecasting outbreak trends.
  \item Reduction in Response Time: Healthcare agencies should receive outbreak warnings to an acceptable time in advance, equal or greater to the minimum allowed time.
\end{itemize}

\subsubsection*{Constraints}

\begin{itemize}
  \item Data Availability: The dataset is restricted to past influenza records and may have incomplete components, erroneous data, or improper data.
  \item Computational Resources: High performing machine learning models may require significant computational power, which can result in higher budget for the project.
  \item Timeliness: There may be an inability to capture real time data because of delays in reporting by different nations.
  \item Privacy \& Regulatory Constraints: Some data items could be sensitive and therefore protected by privacy legislations.
\end{itemize}

\subsubsection*{Organizational and Business Impact}

\begin{itemize}
  \item Positive Impact: Enhanced planning in healthcare systems, reduced economic impact, and improved provision of services.
  \item Negative Impact: If the predictions are imprecise there is going to be severe mismanagement of resources and unnecessary panic.
\end{itemize}

\subsection*{Task 1.3 Determine Data-Mining Goals}

\subsubsection*{Data-Mining Goals}

The purpose of the data mining process is to develop a time series forecasting model for predicting influenza epidemics based on historical monitoring data. The objectives are:
\begin{itemize}
  \item Cleaning and preprocessing the GISRS dataset to remove duplicates and missing values.
  \item Constructing additional time related features (week of the year, month of the year, season) in addition to lag features (cases in previous weeks).
  \item Building and training machine learning models such as Linear Regression, Decision Tree, and random Forest models to forecast the number of influenza cases for the following weeks.
  \item Analyzing models' performance with Mean Absolute Error (MAE) and Mean Squared Error (MSE).
  \item Deploying/running the final model to generate real-time outbreak predictions.
\end{itemize}

\subsubsection*{Data-Mining Success Criteria}

\begin{itemize}
  \item Model Accuracy: Accuracy of the predictive model should attain Mean Absolute Error (MAE) of 5\% or less from the reported received values.
  \item Comparative Performance: The designed model should be more accurate than a simple benchmark model (for example: moving average). The improvement should be no less than 15\%.
  \item Operational Usability: The model output should be useful to public health stakeholders for making decisions.
\end{itemize}
